\newpage
\section{Fase de Pruebas y Validación (Quality Assurance)}

Tras concluir la fase de implementación, el sistema fue sometido a un riguroso proceso de validación empírica (\textit{Quality Assurance}) para garantizar su estabilidad, usabilidad y seguridad antes del despliegue final. Esta fase no solo sirvió para certificar el cumplimiento de la matriz de requisitos, sino también para detectar casos límite (\textit{edge-cases}) y refinar la lógica de negocio mediante ciclos rápidos de \textit{bugfixing}.

\subsection{Validación de Requisitos Funcionales}

Las pruebas funcionales se ejecutaron simulando el comportamiento de un usuario final, abarcando tanto los flujos principales como los alternativos (introducción de datos erróneos o navegación forzada) definidos en los Casos de Uso.

\subsubsection{Módulos 1 y 5: Gestión de Usuarios y Preferencias (CU1-CU5)}
En el módulo de autenticación, se verificó el correcto rechazo de correos duplicados y el manejo de credenciales inválidas. Durante la validación del CU4 (Cerrar Sesión) y CU5 (Eliminar Cuenta), se detectó que los tokens JWT caducados o residuales en el almacenamiento local podían provocar bloqueos en la interfaz (pantallas de carga infinitas) si el servidor reiniciaba la base de datos. Se solucionó implementando un bloque de control en el \texttt{AuthContext} que, ante cualquier fallo de validación del token al arrancar la aplicación, limpia proactivamente el almacenamiento y redirige al usuario a la pantalla pública. 

Asimismo, se validó la persistencia de las preferencias del sistema. Al cambiar el color de acento, la inyección dinámica de variables CSS en la raíz del documento (\texttt{:root}) demostró cambiar el aspecto de toda la plataforma instantáneamente, manteniendo el estado tras recargar el navegador.

\subsubsection{Módulos 2 y 3: Registros, Historial y Borrado Lógico (CU6-CU10)}
Las pruebas del núcleo de la aplicación confirmaron la correcta restricción temporal, impidiendo al usuario crear registros en fechas futuras. El algoritmo de cálculo de rachas demostró sumar los días consecutivos correctamente y reiniciar el contador a 1 cuando se detecta un lapso superior a 24 horas sin actividad.

Uno de los retos técnicos más interesantes surgió al probar el CU8 (Gestionar Etiquetas). Se implementó un sistema de \textbf{borrado lógico} para evitar la pérdida de integridad referencial: cuando un usuario elimina una etiqueta, el sistema cambia su estado a inactivo en lugar de borrar el documento de MongoDB. 

Durante las pruebas del historial, se detectó que los registros pasados que contenían etiquetas inactivas las mostraban como ``Desconocidas''. Esto ocurría porque el cliente solo solicitaba las etiquetas activas al servidor. Para solucionarlo preservando las reglas de negocio, se dividió la lógica en el backend: se habilitó un \textit{endpoint} para el formulario de creación que devuelve exclusivamente las etiquetas activas, y un \textit{endpoint} maestro para el historial que recupera el catálogo completo, garantizando que los registros antiguos se rendericen perfectamente sin permitir que las etiquetas borradas vuelvan a utilizarse.

\subsubsection{Módulo 4: Análisis de Datos y Estados Vacíos (CU11-CU13)}
En el módulo estadístico, se comprobó que los gráficos asimilan correctamente la ausencia de datos. Si el usuario es nuevo, el sistema renderiza ``Empty States'' (pantallas de estado vacío) amigables que invitan a la acción, en lugar de mostrar gráficos en blanco o errores de consola. En el gráfico de donas, el algoritmo heurístico demostró ordenar siempre las categorías de peor a mejor, asociando de forma determinista el color semántico correspondiente a cada porción, facilitando una lectura analítica predecible.

\subsection{Validación de Requisitos No Funcionales (RNF)}

La comprobación de los atributos de calidad del sistema exigió el uso de herramientas de monitorización y de red integradas en el navegador, arrojando resultados sumamente positivos para la viabilidad de la plataforma.

\subsubsection{Rendimiento y Tolerancia a Fallos (Bloque 2)}
Mediante el panel de red (\textit{Network}) de las herramientas de desarrollador, se midió el tiempo de respuesta del servidor (RNF2.1). Las peticiones a la API REST (como la obtención del historial o la recuperación del perfil) promediaron tiempos de respuesta de entre \textbf{13 y 25 milisegundos}, superando holgadamente el límite establecido de 2 segundos. El renderizado del DOM se ejecutó por debajo del límite de 3 segundos (RNF2.2), un tiempo que se verá aún más reducido al compilar y minificar los archivos en el paso a Producción.

Para validar el RNF2.4 (comportamiento ante falta de respuesta), se forzó la caída controlada del contenedor de la base de datos MongoDB durante la navegación. Inicialmente, las peticiones HTTP del cliente se quedaban esperando indefinidamente. Para otorgar robustez a la arquitectura, se implementó un \textbf{Interceptor de Axios Global} en el frontend y se limitó el \textit{Timeout} del cliente Java a 3 segundos. Tras esta mejora, si la base de datos o el servidor dejan de responder, el interceptor captura la excepción \texttt{ECONNABORTED} de forma centralizada y lanza de inmediato una notificación visual, evitando el bloqueo de la pantalla y manteniendo al usuario informado.

\subsubsection{Seguridad, Usabilidad y Portabilidad (Bloques 1, 3 y 4)}
A nivel de usabilidad (RNF3.1), se verificó el comportamiento responsivo mediante emuladores de dispositivos móviles. El menú lateral \textit{desktop} colapsa en una barra de navegación inferior (\textit{Bottom Navigation Bar}) ergonómica sin desbordar los componentes de la cuadrícula. Además, se detectó un problema de renderizado originado por el motor WebKit, el cual provocaba que el texto de las notificaciones flotantes (\textit{Toasts}) se mostrara borroso durante las animaciones al posicionarse en subpíxeles decimales. Esto se solucionó aplicando la propiedad \texttt{transform: translateZ(0)} y \texttt{backface-visibility: hidden} en el CSS global, forzando la aceleración por \textit{hardware} y garantizando la nitidez exigida (RNF3.3).

En cuanto a la seguridad, se demostró mediante inspección del almacenamiento local que la eliminación manual del token JWT expulsa automáticamente al usuario a las rutas públicas (RNF4.2). Además, las consultas en la capa repositorio de Spring Data garantizan el aislamiento de los datos (RNF4.3) concatenando invariablemente el ID extraído del contexto de seguridad de la sesión, haciendo tecnológicamente imposible que un usuario acceda a los registros de terceros.

Por último, la portabilidad (RNF1.2 y RNF1.3) quedó demostrada al ejecutar el sistema de forma idéntica, sin necesidad de instalación de \textit{software} adicional, en los motores de navegación Chromium (Google Chrome) y Gecko (Mozilla Firefox).